\section{引言}
\sourcepage{1--2}
拓扑学的一个基本问题是：对于给定的两个空间，判定它们是否同胚。要证明两个空间同胚，需要构造从一个空间到另一个空间的连续双射，并且其逆映射也连续。要证明两个空间不同胚，则必须证明这样的映射不存在。这通常更为困难。常见的方法是找出某种拓扑性质，即在同胚下保持不变的性质，使得一个空间具有这种性质，而另一个空间没有。例如，$\R^2$ 中的闭单位圆盘不可能与平面 $\R^2$ 同胚，因为闭圆盘是紧致的，而平面不是。实直线 $\R$ 也不可能与 $\R^2$ 同胚，因为从 $\R$ 中删去一个点后，剩余空间不连通，而从 $\R^2$ 中删去一个点则不会如此。

在处理同胚问题时，这些初等性质的作用相当有限。例如，要在同胚意义下对所有紧致曲面进行分类，就需要比上述性质更精细的拓扑不变量。一般地，当 $n\ne m$ 时，要证明 $\R^n$ 与 $\R^m$ 不同胚，也需要这样的不变量。

代数拓扑起源于 Henri Poincaré、Enrico Betti 等数学家构造这类拓扑不变量的尝试。Poincaré 引入了一种群，称为拓扑空间的\textbf{基本群}（fundamental group）；根据其定义，它就是一个拓扑不变量。人们可以相当容易地证明，球面、环面、Klein 瓶等一些熟悉空间的基本群彼此不同，因此这些空间不可能同胚［见 [Mun00] 第 60 节］。事实上，可以利用基本群对所有紧致曲面进行分类［见 [Mas91] 第 4 章第 5 节］。

代数学家 Emmy Noether 推动了研究重点从空间的数值不变量转向空间的同调群。这是因为，在 20 世纪 20 年代中期，Paul Alexandroff 和 Heinz Hopf 曾到 Göttingen 访问，当时 Noether 在该校数学系任职。她认识到，Alexandroff 和 Hopf 所讨论的 Betti 数与挠系数，实际上是与阿贝尔群相联系的不变量；具体说来，就是与同调群相联系的不变量。Hopf 在 1928 年推广 Euler--Poincaré 公式的论文 [Hop28] 的引言中，明确肯定了 Noether 的影响。早期，同调群带来的一个困难在于：它们并非显然是拓扑不变量。一旦证明了这一点，就可以利用这些群来处理同胚问题。

我们将从研究\textbf{同调群}（homology group）开始讨论代数拓扑。随后讨论其他拓扑不变量，例如上同调群和上同调环。我们还将简要介绍同伦群。

定义同调群有几种不同的方法，对于足够良好的空间，这些方法给出相同的结果。我们将详细讨论其中两种：单纯同调群和奇异同调群。首先讨论历史上最先出现的\textbf{单纯同调群}（simplicial homology group）。无论在概念上还是在计算上，它们都很具体、直观。然而，它们只对特别良好的空间，也就是多面体，才有定义；而且，证明它们的拓扑不变性需要付出相当的努力。之后，我们将讨论\textbf{奇异同调群}（singular homology group），它们是作为单纯同调群的推广而引入的。奇异同调群对任意空间都有定义，并且从定义就能直接看出它们是拓扑不变量。此外，在理论研究中，它们比单纯同调群方便得多。虽然奇异同调群不如单纯同调群容易计算，但在两者都有定义时，它们是一致的。

对任意空间定义同调群的第三种方法出自 Eduard Čech。Čech 同调论仍不完全令人满意，但 Čech 上同调论既重要又有用；我们将在本书接近结尾时讨论它。
